% New English translation from checked cote 25, site 98 (PDF page 99).
% AI-translated by GPT-6 Astra (OpenAI, ChatGPT Pro mode).
\scanpage{25/98}
\mspage{1}
\hand{New EGA V}
\var{\inserted{Invertible sheaves and divisors on projective and multiprojective bundles. Linear systems of divisors. \var{\deleted{families of}}{}{deletion}{ega25_25_098-v0002}}}{Invertible sheaves and divisors on projective and multiprojective bundles. Linear systems of divisors. \var{\deleted{families of}}{}{deletion}{ega25_25_098-v0002}}{insertion}{ega25_25_098-v0001}
% NOTE: the reviewed source records two long diagonal pen strokes; their function does not establish cancellation of the entire argument.

IV 2\var{\over{0}{1}}{1}{writing over}{ega25_25_098-v0003} \var{\deleted{\uline{Complements on invertible sheaves}, \uline{the} \uline{divisors of}}}{}{deletion}{ega25_25_098-v0004}
\hand{\unsure{+|t}}\var{\deleted{\uline{Cartier}, \uline{the linear systems of divisors}.}}{}{deletion}{ega25_25_098-v0005}

1. \uline{Determination of invertible sheaves on projective bundles}.
\uline{Application} \uline{to the \var{\deleted{group}}{}{deletion}{ega25_25_098-v0007} of}\hand{/}\uline{automorphisms of a projective bundle}.

\uline{Theorem} 1.1. Let S be a prescheme, \underline{E} a loc free Module on \var{\xout{S which}}{}{typewriter strike-out}{ega25_25_098-v0008} of
\hand{finite} rank $\hand{\geqslant 2}$ at every point,
\var{\xout{is zero at no point,}}{}{typewriter strike-out}{ega25_25_098-v0009} P = P(\underline{E}) the projective bundle it defines. Then
\var{\margin{\unsure{*|+}}}{\placedmargin{\unsure{*|+}}}{marginal note placed}{ega25_25_098-v0010} for every invertible Module \underline{L} on P, we can find a family
\var{\xout{\illegible{1}}}{}{typewriter strike-out}{ega25_25_098-v0011} $(S_n)_{n\hand{\in}\underline{Z}}$ of disjoint opens of S covering S, indexed by \underline{Z}, and
\var{\xout{for every n$\in$\underline{Z}}}{}{typewriter strike-out}{ega25_25_098-v0012} an invertible Module \var{\xout{$\underline{L}_n$}}{}{typewriter strike-out}{ega25_25_098-v0013}\var{\inserted[typed]{$\underline{M}$}}{$\underline{M}$}{insertion}{ega25_25_098-v0014} on S\var{\xout{$_n$}}{}{typewriter strike-out}{ega25_25_098-v0015}, in such a way that
\var{\deleted{\var{\xout{\unsure{$\underline{L}|p^{-1}(S_n)$ is isomorphic to $p_n^{-1}(\underline{L}_n)\otimes\underline{O}$ the restriction of}}}{}{typewriter strike-out}{ega25_25_098-v0017}}}{}{deletion}{ega25_25_098-v0016} the restriction
of \underline{L} to $f^{-1}(S_n)$ is isomorphic to that of $\underline{M}\otimes_{\underline{O}_S}\underline{O}_P(1) = f^*(\underline{M})(n)$.
Moreover, the \var{\xout{family}}{}{typewriter strike-out}{ega25_25_098-v0018} family of the $S_n$ \var{\xout{and the class}}{}{typewriter strike-out}{ega25_25_098-v0019} is uniquely
determined by its conditions, as is \underline{M} up to unique isomorphism.

\uline{Remark} 1.2. When we make no hypothesis on the rank of \underline{E}, S
decomposes canonically as \var{\inserted{the}}{the }{insertion}{ega25_25_098-v0020}\var{\inserted[typed]{prescheme}}{prescheme }{insertion}{ega25_25_098-v0021}sum of three opens $S^0,S^1,S^2$ such that on
these, the rank of \underline{E} is resp. 0, 1, and $\hand{\geqslant} 2$. Then the determination of
invertible Modules on P reduces to that of the $f^{-1}(S^{\unsure{i|n}}_{\unsure{a|n}})$ for i=0,1,2.
The case i=2 can be dealt with by 1.1., on the other hand $f^{-1}(S^1)$ is $S^1$-isomorphic
to $S^1$ so its Picard group is precisely Pic($S^1$), finally
$f^{-1}(S^0)$ is empty so its Picard group is zero.

\uline{Corollary} 1.3. Under the conditions of 1.1. suppose S connected\var{\inserted[typed]{nonempty}}{ nonempty}{insertion}{ega25_25_098-v0022}. Then
every invertible Module \underline{L} on P is isomorphic to a Module of the form
$f^*(\underline{M})(n)$, where $n\hand{\in}\underline{Z}$ and \underline{M} is an invertible Module on S. Moreover n is
uniquely determined, and \underline{M} is determined up to unique isom, by the datum
of \underline{L}.\var{\margin{\ensuremath{\swarrow}}}{\placedmargin{\ensuremath{\swarrow}}}{marginal note placed}{ega25_25_098-v0023}

Another way of formulating this corollary is the following. Consider
the natural homomorphisms Pic(S)$\hand{\longrightarrow}$Pic(\unsure{P|R|E}), \underline{Z} $\hand{\longrightarrow}$Pic(P), the first\var{\margin{\ensuremath{\swarrow}}}{\placedmargin{\ensuremath{\swarrow}}}{marginal note placed}{ega25_25_098-v0024}
\tnote{25-098-1}{The incomplete wording after the deleted group in the heading, and the missing object after on in the theorem's corrected first sentence, are retained. The source's phrase its conditions and the twist (1) beside the final (n) in the displayed-in-text formula are also retained. The canceled stray x after Application has no English spelling effect.}
\var{\margin{\unsure{$\backslash$|1}}}{\placedmargin{\unsure{$\backslash$|1}}}{marginal note placed}{ega25_25_098-v0025}\var{\margin{98}}{\placedmargin{98}}{marginal note placed}{ega25_25_098-v0026}
